\slajdid{08-s035}
\begin{frame}[label=08-s035]{Dinamičke karakteristike logičkog kola}
\gusttekst\setlength{\parskip}{2pt}
\begin{columns}[c,onlytextwidth]\column{.35\textwidth}Neke definicije\par\medskip\begin{tikzpicture}[circuit logic US,DEoznaka,baseline=(g.center)]\node[not gate US,draw](g){};\draw(g.input)--++(-.35,0)node[left]{$a_i(t)$};\draw(g.output)--++(.35,0)node[right]{$b_o(t)$};\end{tikzpicture}

\column{.63\textwidth}\centering\begin{tikzpicture}[x=.9cm,y=.85cm,font=\sitneoznake,>=Latex]
\draw[->](-1.65,0)--(5.25,0)node[below]{$t$};\draw[->](0,-.1)--(0,2.1)node[left]{$a_i(t)$};
\node[above]at(-1.1,1.2){$\approx V_{OH}$};\node[above]at(.8,.25){$\approx V_{OL}$};
\draw(-1.5,1.2)--(-.9,1.2)..controls(-.6,1.2)and(-.2,0.25)..(.1,0.25)--(1.4,0.25)..controls(1.7,0.25)and(2.1,1.2)..(2.4,1.2)--(3.4,1.2)..controls(3.7,1.2)and(4.1,0.25)..(4.4,0.25)--(5,0.25);
\node[align=center](r)at(1.6,2){usponska, uzlazna,\\ivica signala};\draw[->](r.south)--(1.9,1.05);
\node[align=center](f)at(4.4,2){nisponska, silazna,\\ivica signala};\draw[->](f.south)--(3.92,.85);
\end{tikzpicture}
\end{columns}
Promena signala je između niskog nivoa, $V_{OL}$, i visokog nivoa, $V_{OH}$. Amplituda promene je $A=V_{OH}-V_{OL}$.\par
Sledeće dve tačke od interesa su $V_{OL}+10\%A$ i $V_{OL}+90\%A$.\par
\begin{columns}[c,onlytextwidth]\column{.52\textwidth}\centering\begin{tikzpicture}[x=.9cm,y=1.0cm,font=\sitneoznake,>=Latex]
\draw[->](-1.65,-.2)--(5.25,-.2)node[below]{$t$};\draw[->](0,-.3)--(0,2.1)node[left]{$a_i(t)$};
\node[above]at(-1.1,1.2){$\approx V_{OH}$};\node[below]at(.8,.25){$\approx V_{OL}$};
\draw(-1.5,1.2)--(-.9,1.2)..controls(-.6,1.2)and(-.2,0.25)..(.1,0.25)--(1.4,0.25)..controls(1.7,0.25)and(2.1,1.2)..(2.4,1.2)--(3.4,1.2)..controls(3.7,1.2)and(4.1,0.25)..(4.4,0.25)--(5,0.25);

\draw[dashed](-.12,.345)--(4.214,.345);\node[above right]at(0,.345){$V_{OL}+10\%A$};\draw[dashed](-.12,1.105)--(3.586,1.105);\node[above right]at(0,1.105){$V_{OL}+90\%A$};
\foreach\x/\y in{1.586/.345,2.214/1.105,3.586/1.105,4.214/.345}{\draw[dashed](\x,-.52)--(\x,\y);}
\draw[<->](1.586,-.38)--(2.214,-.38)node[midway,below]{$t_r$};\draw[<->](3.586,-.38)--(4.214,-.38)node[midway,below]{$t_f$};
\node[below]at(2,-.85){Principski $t_r$ različito od $t_f$};
\end{tikzpicture}

\column{.46\textwidth}
Vreme koje usponska ivica provede između ove dve tačke naziva se vremenom uspona, trajanjem usponske ivice, rise time i obeležava se sa $t_r$.\par\medskip
Vreme koje nisponska ivica provede između ove dve tačke naziva se vremenom pada, trajanjem nisponske ivice, fall time i obeležava se sa $t_f$.
\end{columns}
\end{frame}
